\documentclass[11pt]{article}
\usepackage[margin=1in]{geometry}
\usepackage{booktabs}
\usepackage{rotating}
\usepackage{pdflscape}
\usepackage{longtable}

\title{Rural Labor Markets and Small-Business Formation in Texas:\\
A County-Level Assessment}
\author{Office of Rural Economic Research}
\date{July 2026}

\begin{document}
\maketitle

\section{Introduction}

Rural Texas has long been treated as a residual category in state economic
planning: the counties that remain once the metropolitan triangle of Dallas,
Houston, San Antonio, and Austin is set aside. Yet these counties hold a
disproportionate share of the state's land, water, and energy infrastructure,
and their labor markets behave in ways that statewide averages obscure. This
report assembles county-level indicators on population, labor force
participation, household income, broadband availability, business
establishment counts, and recent employment change to characterize how rural
labor markets have evolved since the pandemic-era disruptions.

The analysis draws on simulated administrative records constructed for
methodological testing. No figure in this report should be cited as an
official statistic; the purpose is to demonstrate a reporting format that
state agencies could adopt for genuine data releases.

\section{Regional patterns in the western districts}

The clearest divergence appears in the western planning districts, where
energy-sector volatility drives large swings in both employment and
establishment counts. Table~\ref{tab:western} presents the full indicator
panel for a selection of western district counties. Two patterns stand out.
First, labor force participation varies far more across neighboring counties
than conventional regional summaries suggest. Second, establishment growth is
only weakly correlated with population size: several of the smallest counties
posted the strongest gains, consistent with a low-base effect rather than
genuine structural change.

\begin{sidewaystable}
\centering
\caption{Selected western district indicators, simulated 2025 panel}
\label{tab:western}
\begin{tabular}{lrrrrrrr}
\toprule
County & Population & LFP (\%) & Median HH inc.\ (\$) & Broadband gap (\%) & Establishments & Empl.\ change (\%) & District code \\
\midrule
Ector Bend    & 61,447 & 58.3 & 52,914 & 12.7 & 2,183 &  3.9 & W-14 \\
Pecos Draw    &  8,731 & 44.9 & 47,208 & 21.4 &   412 & 11.2 & W-14 \\
Sierra Blanca &  4,109 & 39.1 & 43,551 & 28.8 &   187 & -2.1 & W-15 \\
Toyah Flats   & 12,586 & 51.7 & 55,873 & 17.2 &   639 &  6.8 & W-15 \\
Salt Basin    & 27,914 & 63.4 & 61,207 &  9.3 & 1,046 &  1.4 & W-16 \\
Quito Mesa    &  3,358 & 41.8 & 49,776 & 24.9 &   141 &  8.7 & W-16 \\
Alamito Creek & 19,462 & 55.2 & 58,341 & 14.6 &   823 & -0.8 & W-17 \\
\bottomrule
\end{tabular}
\end{sidewaystable}

\section{Establishment dynamics}

A second question concerns whether small-business formation in rural counties
is keeping pace with the state as a whole. Table~\ref{tab:formation} reports
establishment birth and death rates alongside net formation for a set of
comparison groups defined by district and by distance to the nearest
metropolitan area. The headline finding is that net formation in remote
counties is positive but fragile: birth rates are respectable, but death rates
run several points higher than in counties adjacent to a metro, so the net
figure depends heavily on a handful of surviving firms in each cohort.

\begin{landscape}
\begin{table}
\centering
\caption{Establishment birth and death rates by comparison group, simulated 2023--2025 cohorts}
\label{tab:formation}
\begin{tabular}{lrrrrrrrr}
\toprule
Comparison group & Births 2023 & Deaths 2023 & Births 2024 & Deaths 2024 & Births 2025 & Deaths 2025 & Net rate (\%) & Survival at 24 mo.\ (\%) \\
\midrule
Metro-adjacent, ag-dominant      & 1,872 & 1,431 & 1,944 & 1,502 & 2,036 & 1,489 &  4.6 & 71.8 \\
Metro-adjacent, energy-dominant  & 2,215 & 1,988 & 2,301 & 2,144 & 2,187 & 2,079 &  1.9 & 64.2 \\
Remote, ag-dominant              &   934 &   861 &   917 &   889 &   961 &   902 &  1.1 & 68.4 \\
Remote, energy-dominant          & 1,148 & 1,207 & 1,233 & 1,164 & 1,296 & 1,171 &  0.7 & 59.6 \\
Remote, trade-corridor           &   688 &   594 &   731 &   617 &   779 &   633 &  3.8 & 73.1 \\
Border, mixed economy            & 1,509 & 1,342 & 1,587 & 1,398 & 1,644 & 1,417 &  3.2 & 69.9 \\
\bottomrule
\end{tabular}
\end{table}
\end{landscape}

\section{The full county panel}

For completeness, Table~\ref{tab:panel} lists the complete simulated panel
for all counties in the study frame. Readers interested in a particular
county can locate it there; the columns follow the same definitions used in
the preceding sections. Because the panel is wide, it is set in landscape
orientation and continues across multiple pages.

Three cautions apply when reading the panel. Small-county rates are noisy:
a single plant opening or closure can move the employment-change column by
several points. Broadband gap figures reflect advertised rather than
delivered service. And the establishment counts include nonemployer firms,
which inflates totals in counties with substantial farm and ranch
proprietorships.

\begin{landscape}
\begin{longtable}{lrrrrrr}
\caption{Complete simulated county panel: population, labor force participation, income, broadband, establishments, and employment change}
\label{tab:panel} \\
\toprule
County & Population & LFP (\%) & Median HH inc.\ (\$) & Broadband gap (\%) & Establishments & Empl.\ change (\%) \\
\midrule
\endfirsthead
\multicolumn{7}{l}{\small\emph{Table \ref{tab:panel}, continued}} \\
\toprule
County & Population & LFP (\%) & Median HH inc.\ (\$) & Broadband gap (\%) & Establishments & Empl.\ change (\%) \\
\midrule
\endhead
\midrule
\multicolumn{7}{r}{\small\emph{Continued on next page}} \\
\endfoot
\bottomrule
\endlastfoot
Nacogdoches & 53,827 & 62.5 & 43,967 & 18.6 & 1,307 & 2.6 \\
Zavala & 96,887 & 74.4 & 47,842 & 14.8 & 2,391 & 5.1 \\
Kleberg & 26,948 & 32.0 & 40,466 & 8.4 & 1,503 & 13.8 \\
Ochiltree & 94,746 & 65.9 & 54,209 & 24.4 & 1,598 & 2.2 \\
Wharton & 51,766 & 41.3 & 51,013 & 22.4 & 2,913 & 7.6 \\
Bosque & 41,257 & 59.6 & 50,196 & 21.3 & 2,444 & -0.4 \\
Refugio & 27,649 & 58.0 & 78,879 & 13.2 & 3,874 & -6.7 \\
Terrell & 93,163 & 32.5 & 73,585 & 24.1 & 3,382 & 4.1 \\
Yoakum & 59,769 & 59.3 & 38,849 & 14.3 & 1,018 & 9.8 \\
Menard & 68,326 & 61.0 & 76,834 & 15.2 & 4,048 & -9.5 \\
Lampasas & 30,561 & 60.4 & 60,264 & 24.0 & 917 & -4.4 \\
Gaines & 9,149 & 41.4 & 66,970 & 4.7 & 3,802 & -5.2 \\
Motley & 18,999 & 43.3 & 71,271 & 17.9 & 2,252 & -3.6 \\
Sabine & 10,536 & 74.8 & 57,617 & 17.3 & 2,700 & -3.7 \\
Duval & 59,214 & 36.6 & 43,405 & 15.3 & 2,857 & -0.9 \\
Crockett & 21,461 & 69.9 & 70,314 & 26.3 & 295 & 4.3 \\
Hansford & 90,175 & 42.7 & 77,954 & 23.6 & 2,708 & 8.2 \\
Jim Hogg & 73,958 & 32.7 & 57,653 & 5.9 & 2,845 & 6.5 \\
La Salle & 76,655 & 42.6 & 49,153 & 10.9 & 2,232 & 5.3 \\
Loving & 29,529 & 49.0 & 65,350 & 11.0 & 3,918 & -4.7 \\
Upshur & 49,031 & 72.9 & 43,652 & 22.9 & 203 & -3.4 \\
Winkler & 54,505 & 48.0 & 75,568 & 22.7 & 1,334 & 3.6 \\
Coke & 42,633 & 32.6 & 67,791 & 4.7 & 2,061 & 12.5 \\
Delta & 90,649 & 73.3 & 43,012 & 7.5 & 3,591 & 0.0 \\
Foard & 74,044 & 76.6 & 61,065 & 18.0 & 2,106 & 8.2 \\
Glasscock & 2,638 & 72.7 & 62,841 & 23.6 & 1,845 & -4.6 \\
Hemphill & 23,411 & 52.1 & 79,450 & 15.5 & 1,679 & 11.9 \\
Irion & 81,507 & 45.0 & 60,519 & 19.9 & 1,352 & -3.7 \\
Jeff Davis & 95,654 & 62.6 & 51,039 & 5.0 & 801 & -3.0 \\
Kenedy & 72,134 & 60.2 & 80,064 & 23.5 & 768 & 5.1 \\
Lipscomb & 45,214 & 57.5 & 80,862 & 26.0 & 185 & 5.4 \\
Mills & 32,296 & 34.8 & 61,132 & 16.7 & 2,363 & 13.6 \\
Newton & 92,903 & 71.6 & 57,213 & 20.3 & 164 & -6.6 \\
Oldham & 90,042 & 41.3 & 41,329 & 22.9 & 2,179 & 5.4 \\
Presidio & 62,861 & 54.4 & 40,383 & 4.2 & 958 & -4.6 \\
Rains & 62,704 & 46.1 & 56,820 & 6.0 & 3,785 & -1.3 \\
Real & 6,528 & 40.3 & 41,511 & 8.9 & 1,752 & 2.3 \\
Roberts & 29,789 & 39.2 & 80,752 & 21.9 & 3,010 & -5.7 \\
Runnels & 18,394 & 45.7 & 70,473 & 14.8 & 689 & -5.7 \\
San Saba & 84,946 & 59.7 & 62,213 & 24.7 & 365 & 9.2 \\
Schleicher & 18,743 & 53.3 & 52,117 & 7.5 & 1,797 & -2.6 \\
Shackelford & 6,171 & 41.4 & 46,522 & 10.2 & 2,011 & 8.9 \\
Somervell & 50,021 & 42.5 & 61,677 & 5.8 & 2,974 & -4.8 \\
Sterling & 33,040 & 53.2 & 65,448 & 21.3 & 2,307 & 13.7 \\
Stonewall & 47,967 & 44.4 & 58,542 & 9.6 & 123 & 13.6 \\
Sutton & 89,708 & 34.9 & 56,396 & 7.8 & 2,360 & 12.6 \\
Swisher & 47,406 & 43.6 & 70,196 & 8.5 & 2,800 & -8.0 \\
Throckmorton & 57,642 & 42.6 & 48,817 & 17.9 & 2,590 & 13.7 \\
Trinity & 11,894 & 70.4 & 75,588 & 25.8 & 1,134 & -9.2 \\
Upton & 42,476 & 56.4 & 70,083 & 21.6 & 1,880 & 9.4 \\
Wheeler & 76,443 & 49.6 & 64,563 & 10.1 & 899 & -3.8 \\
Wilbarger & 76,441 & 37.9 & 48,779 & 11.8 & 1,127 & -5.9 \\
Zapata & 90,824 & 77.6 & 45,206 & 20.1 & 3,762 & 13.0 \\
Armstrong & 70,716 & 61.0 & 79,547 & 26.6 & 270 & -6.0 \\
Bailey & 55,116 & 67.0 & 73,847 & 7.4 & 1,137 & -5.1 \\
Baylor & 32,142 & 69.8 & 64,998 & 15.3 & 1,970 & -1.2 \\
Borden & 13,348 & 34.4 & 45,053 & 19.3 & 1,303 & 4.2 \\
Briscoe & 79,490 & 73.9 & 44,359 & 13.3 & 3,801 & 4.5 \\
Castro & 17,614 & 76.1 & 76,043 & 8.0 & 3,568 & -2.0 \\
Childress & 3,422 & 70.8 & 59,065 & 11.3 & 342 & 12.8 \\
Cochran & 13,040 & 45.0 & 58,498 & 10.4 & 1,105 & 2.7 \\
Collingsworth & 52,277 & 77.5 & 66,992 & 20.6 & 488 & 2.5 \\
\end{longtable}
\end{landscape}

\section{Conclusion}

The format demonstrated here---compact indicator tables for regional
highlights (Table~\ref{tab:western}), grouped comparisons for dynamics
(Table~\ref{tab:formation}), and a full landscape panel for reference
(Table~\ref{tab:panel})---balances readability with completeness. Agencies
adopting it for real data should pair each table with the microdata vintage
and suppression rules used to construct it, so that county officials can
reconcile published figures against their own records.

\end{document}
